% Annotations on Evan Chen's Infinitely Large Napkin
% Andre Ruiz Loera

\documentclass[11pt, letterpaper]{article}

\usepackage[margin=1.2in]{geometry}
\usepackage{amsmath, amssymb, amsthm}
\usepackage{mathtools}
\usepackage{hyperref}
\usepackage{enumitem}
\usepackage{xcolor}

% Paragraph spacing without parskip package
\setlength{\parskip}{0.6em}
\setlength{\parindent}{0pt}

% Theorem environments
\newtheorem{theorem}{Theorem}[section]
\newtheorem{lemma}[theorem]{Lemma}
\newtheorem{proposition}[theorem]{Proposition}
\newtheorem{corollary}[theorem]{Corollary}
\theoremstyle{definition}
\newtheorem{definition}[theorem]{Definition}
\newtheorem{example}[theorem]{Example}
\theoremstyle{remark}
\newtheorem*{remark}{Remark}
\newtheorem*{note}{Note}

% Common math macros
\newcommand{\R}{\mathbb{R}}
\newcommand{\N}{\mathbb{N}}
\newcommand{\Z}{\mathbb{Z}}
\newcommand{\Q}{\mathbb{Q}}
\newcommand{\C}{\mathbb{C}}
\newcommand{\E}{\mathbb{E}}
\newcommand{\Prob}{\mathbb{P}}
\newcommand{\norm}[1]{\left\lVert#1\right\rVert}
\newcommand{\abs}[1]{\left|#1\right|}
\newcommand{\inner}[2]{\langle #1,\, #2 \rangle}

\title{Annotations on \textit{An Infinitely Large Napkin} \\ \large Evan Chen}
\author{Andre Ruiz Loera}
\date{May 2026}

\begin{document}
\maketitle

\begin{abstract}
  % Fill this in — why are you reading the Napkin and what do you want to get out of it?
  The Napkin is Evan Chen's ambitious attempt to expose a motivated undergraduate to
  the breadth of modern mathematics in a single document: algebra, topology, analysis,
  combinatorics, and more, each developed from scratch but pushed quickly toward
  non-trivial results. I am reading it to \textbf{[fill in your reason]}.
\end{abstract}

\tableofcontents
\newpage

% ══════════════════════════════════════════════════════════════
\section{Why the Napkin?}
% A few sentences on what drew you to this text.


% ══════════════════════════════════════════════════════════════
\section{Part I — Starting Out}

\subsection{Chapter 1 — Groups}

\begin{definition}[Group]
  A \emph{group} is a set $G$ together with a binary operation $\cdot : G \times G \to G$
  satisfying: % fill in axioms as you annotate
\end{definition}

% Add your annotations, proofs, examples below each subsection.


% ══════════════════════════════════════════════════════════════
\section{Part II — Abstract Algebra Continued}

\subsection{Chapter — }


% ══════════════════════════════════════════════════════════════
\section{Part III — Linear Algebra}

\subsection{Chapter — }


% ══════════════════════════════════════════════════════════════
\section{Part IV — Complex Analysis}

\subsection{Chapter — }


% ══════════════════════════════════════════════════════════════
\section{Part V — Topology}

\subsection{Chapter — }


% ══════════════════════════════════════════════════════════════
\section{Problems \& Exercises}
% Work through Napkin exercises here.

\subsection*{Exercise X.Y}
\textit{State the problem here.}

\begin{proof}
  % Your solution.
\end{proof}


% ══════════════════════════════════════════════════════════════
\section{Running List of Key Results}
% Theorems and lemmas worth remembering across chapters.

\begin{theorem}[Name of theorem]
  Statement.
\end{theorem}

\begin{proof}
  % Proof or sketch.
\end{proof}

\end{document}
